\chapter{Dual Frames}
\label{ch:dual-frames}
\fancyhead[R]{\small\textit{Chapter 5: Dual Frames}}
\section{Overview: From One Reconstruction Formula to a Family}
\label{sec:ch5-overview}

Chapter~3 derived a reconstruction formula almost in passing, while
studying the frame operator: for any frame $\{f_i\}_{i=1}^m$ with frame
operator $S$,
\[
  x = \sum_{i=1}^m \inner{x}{f_i}\,(S^{-1}f_i) \qquad \text{for every } x \in V,
\]
and named the vectors $\tilde f_i := S^{-1}f_i$ the \emph{canonical dual
frame}, promising a full treatment in this chapter. Chapter~3 also
observed, in passing, that the synthesis operator $F$ has a nontrivial
null space whenever the frame is redundant ($m>n$), and promised that
this null space would turn out to parametrize \emph{every} valid dual,
not just the canonical one.

This chapter delivers on both promises. We will see that the canonical
dual is only one member---admittedly a distinguished one---of an entire
\emph{affine family} of dual frames, all of which reconstruct every
vector in $V$ exactly. We characterize this family completely
(Theorem~\ref{thm:dual-family}), show that the canonical dual is the
unique \emph{minimal-norm} member of the family
(Theorem~\ref{thm:canonical-minimal-norm}), and close with a numerical
investigation of how the choice of dual affects robustness to noise---
returning, with real machinery in hand, to the very motivation that
opened Chapter~1.

\section{Dual Frames: Definition}
\label{sec:dual-definition}

\begin{definition}[Dual frame]
\label{def:dual-frame}
Let $\{f_i\}_{i=1}^m$ be a frame for $V=\F^n$. A frame $\{g_i\}_{i=1}^m$
for $V$ (indexed by the same index set) is called a \emph{dual frame} of
$\{f_i\}_{i=1}^m$ if
\begin{equation}
\label{eq:dual-condition}
  x = \sum_{i=1}^m \inner{x}{f_i}\, g_i \qquad \text{for every } x \in V.
\end{equation}
\end{definition}

\begin{remark*}
Definition~\ref{def:dual-frame} is not symmetric-looking at first
glance, but it is in fact a symmetric relationship: if $\{g_i\}$ is dual
to $\{f_i\}$ in the sense above, then $\{f_i\}$ is also dual to
$\{g_i\}$, i.e.\ $x=\sum_i\inner{x}{g_i}f_i$ for all $x$ as well. We
prove this symmetry as part of Proposition~\ref{prop:canonical-dual-properties}
below, in the canonical case, and the general case is
Lab Exercise 5.1 (it follows the same pattern).
For now, simply note that ``dual'' is meant in this mutual sense, exactly
as ``$\tilde f_i$ is dual to $f_i$'' is meant to be an if-and-only-if kind
of partnership, not a one-way relationship.
\end{remark*}

In matrix language, if $G$ denotes the synthesis matrix of $\{g_i\}_{i=1}^m$
(the $n\times m$ matrix with $i$-th column $g_i$), then
\eqref{eq:dual-condition} reads $x = G(F^*x)$ for every $x\in V$, i.e.
\begin{equation}
\label{eq:dual-matrix-condition}
  GF^* = I_n.
\end{equation}
This is the condition we will use throughout the chapter: \emph{a
synthesis matrix $G$ defines a dual frame if and only if $GF^*=I_n$}.

\section{The Canonical Dual Frame}
\label{sec:canonical-dual}

\begin{proposition}[The canonical dual is a valid dual frame]
\label{prop:canonical-dual-properties}
Let $\{f_i\}_{i=1}^m$ be a frame for $V$ with frame operator $S$, and
define $\tilde f_i := S^{-1}f_i$ for $i=1,\ldots,m$. Then:
\begin{enumerate}[label=(\roman*)]
  \item $\{\tilde f_i\}_{i=1}^m$ is itself a frame for $V$ (the
        \emph{canonical dual frame}).
  \item Its frame operator is $\tilde S = S^{-1}$.
  \item $\{\tilde f_i\}$ is dual to $\{f_i\}$ in the sense of
        Definition~\ref{def:dual-frame}, and the relationship is
        mutual: $\{f_i\}$ is also dual to $\{\tilde f_i\}$.
  \item The optimal frame bounds of $\{\tilde f_i\}$ are
        $\tilde A = 1/B$ and $\tilde B = 1/A$, where $A,B$ are the
        optimal bounds of $\{f_i\}$.
\end{enumerate}
\end{proposition}

\begin{proof}
Write $\tilde F := S^{-1}F$ for the synthesis matrix of $\{\tilde
f_i\}$ (this is consistent: the $i$-th column of $S^{-1}F$ is
$S^{-1}f_i=\tilde f_i$, since $S^{-1}$ is applied to each column of $F$).

\textbf{(ii)} The frame operator of $\{\tilde f_i\}$ is
\[
  \tilde S = \tilde F \tilde F^* = (S^{-1}F)(S^{-1}F)^*
  = S^{-1}F F^*(S^{-1})^* = S^{-1}SS^{-1} = S^{-1},
\]
using that $S$ (hence $S^{-1}$) is self-adjoint, so $(S^{-1})^*=S^{-1}$.

\textbf{(i)} Since $S$ is invertible and self-adjoint with $S>0$
(Proposition~3.4.2(iii)), $S^{-1}$ is also self-adjoint and positive
definite (its eigenvalues are the reciprocals $1/\lambda_k>0$ of the
eigenvalues of $S$). By part (ii), the frame operator of
$\{\tilde f_i\}$ is $\tilde S = S^{-1}$, which is positive definite;
by Proposition~3.4.2(iii) applied in reverse (positive definite frame
operator $\Leftrightarrow$ the vectors span $V$), $\{\tilde f_i\}$
spans $V$, hence (Theorem~2.2.1) is a frame for $V$.

\textbf{(iii)} We must check $\tilde F F^* = I_n$ (matching
\eqref{eq:dual-matrix-condition} with $G=\tilde F$):
\[
  \tilde F F^* = (S^{-1}F)F^* = S^{-1}(FF^*) = S^{-1}S = I_n.
\]
So $\{\tilde f_i\}$ is dual to $\{f_i\}$. For the converse direction,
we must check $F\tilde F^* = I_n$:
\[
  F\tilde F^* = F(S^{-1}F)^* = F F^*(S^{-1})^* = S S^{-1} = I_n,
\]
again using $(S^{-1})^*=S^{-1}$. So $\{f_i\}$ is also dual to
$\{\tilde f_i\}$, confirming the mutual relationship.

\textbf{(iv)} By Theorem~3.5.1, the optimal bounds of $\{\tilde f_i\}$
are the extreme eigenvalues of $\tilde S = S^{-1}$. If $S$ has
eigenvalues $\lambda_1\ge\cdots\ge\lambda_n$ (so $B=\lambda_1$,
$A=\lambda_n$ for $\{f_i\}$), then $S^{-1}$ has eigenvalues
$1/\lambda_1\le\cdots\le1/\lambda_n$ (matrix inversion preserves
eigenvectors and inverts eigenvalues). The smallest eigenvalue of
$S^{-1}$ is $1/\lambda_1=1/B$ and the largest is $1/\lambda_n=1/A$,
giving $\tilde A=1/B$ and $\tilde B=1/A$.
\end{proof}

\begin{keyidea}
The canonical dual frame $\tilde f_i = S^{-1}f_i$ has frame operator
$\tilde S = S^{-1}$, and its frame bounds are the \emph{reciprocals,
in reversed order}, of the original frame's bounds: $\tilde A=1/B$,
$\tilde B=1/A$. In particular, \textbf{the canonical dual of a tight
frame is again tight}: if $A=B$, then $\tilde A=\tilde B=1/A$. This matches the reconstruction formula from Chapter~3 (Equation 3.5),
where the tight-frame dual was simply $\tilde f_i =
\tfrac1A f_i$ --- a special case of $\tilde f_i = S^{-1}f_i = \tfrac1A f_i$
when $S=AI$.
\end{keyidea}

\begin{examplebox}[title={Example 5.1: Canonical dual of the non-tight frame}]
\label{ex:canonical-dual-h}
Recall $\{h_1,h_2,h_3\}=\{(1,0)^T,(0,1)^T,(1,1)^T\}$ from Chapter~3, with
frame operator $S=\begin{pmatrix}2&1\\1&2\end{pmatrix}$ and bounds
$A=1$, $B=3$. We invert $S$:
\[
  S^{-1} = \frac{1}{\det S}\begin{pmatrix}2&-1\\-1&2\end{pmatrix}
  = \frac{1}{3}\begin{pmatrix}2&-1\\-1&2\end{pmatrix},
\]
using $\det S = 2\cdot2-1\cdot1=3$. The canonical dual vectors are
\[
  \tilde h_1 = S^{-1}h_1 = \tfrac13(2,-1)^T, \quad
  \tilde h_2 = S^{-1}h_2 = \tfrac13(-1,2)^T, \quad
  \tilde h_3 = S^{-1}h_3 = \tfrac13(1,1)^T.
\]
By Proposition~\ref{prop:canonical-dual-properties}(iv), the dual frame
has bounds $\tilde A = 1/3$, $\tilde B = 1/1 = 1$. Notice that
$\{\tilde h_i\}$ is visibly \emph{not} a rescaling of $\{h_i\}$---unlike
the tight-frame case, the canonical dual of a non-tight frame is
genuinely a different set of vectors, not merely the original frame
scaled by a constant.
\end{examplebox}

\section{The Family of All Dual Frames}
\label{sec:dual-family}

The canonical dual is a distinguished dual frame, but---exactly as
Chapter~3 promised---it is not the only one whenever the original frame
is redundant. We now characterize the complete family.

\begin{theorem}[Characterization of all dual frames]
\label{thm:dual-family}
Let $\{f_i\}_{i=1}^m$ be a frame for $V=\F^n$ with synthesis matrix $F$,
and let $\tilde F = S^{-1}F$ be the canonical dual synthesis matrix.
Then $G$ is the synthesis matrix of a dual frame of $\{f_i\}$ if and
only if
\begin{equation}
\label{eq:dual-family}
  G = \tilde F + Z, \qquad \text{where } Z F^* = 0.
\end{equation}
That is, the set of all valid dual synthesis matrices is the affine
subspace $\tilde F + \setdef{Z \in \F^{n\times m}}{ZF^*=0}$.
\end{theorem}

\begin{proof}
($\Leftarrow$) Suppose $G=\tilde F+Z$ with $ZF^*=0$. Then
\[
  GF^* = \tilde F F^* + ZF^* = I_n + 0 = I_n,
\]
using Proposition~\ref{prop:canonical-dual-properties}(iii)
($\tilde FF^*=I_n$). By \eqref{eq:dual-matrix-condition}, $G$ defines a
dual frame.

($\Rightarrow$) Suppose $G$ defines a dual frame, so $GF^*=I_n$. Set
$Z := G - \tilde F$. Then
\[
  ZF^* = GF^* - \tilde FF^* = I_n - I_n = 0,
\]
so $Z$ has the required property, and $G=\tilde F+Z$ by construction.
\end{proof}

\begin{remark*}
This theorem makes precise the promise from Chapter~3's remark on the
synthesis operator: the freedom in choosing a dual frame is exactly
captured by matrices $Z$ satisfying $ZF^*=0$, i.e.\ matrices whose rows
each lie in $\ker F$ (since $(ZF^*)_{ji}=\text{row}_j(Z)\cdot
\overline{f_i}$, and this vanishes for every $i$ exactly when
row$_j(Z) \in (\spn\{f_i\})^\perp$ in the appropriate sense--- more
directly, $ZF^*=0$ says each row of $Z$, viewed as a vector of
coefficients in $\F^m$, lies in $\ker F$, since $F\,(\text{row}_j(Z))^*=
0$ is what $ZF^*=0$ unpacks to row by row). When the frame is a basis
($m=n$), $F$ is square and invertible, so $\ker F=\{0\}$, forcing $Z=0$:
\emph{the dual of a basis is unique}, namely the canonical dual, exactly
as expected from ordinary linear algebra (a basis has one and only one
dual/biorthogonal basis). When $m>n$, $\ker F$ is nontrivial
(dimension $m-n$, by the rank-nullity theorem, since $F$ has rank $n$),
and there is a genuine $(m-n)\times n$-real-dimensional family of dual
frames to choose from.
\end{remark*}

\begin{keyidea}
\textbf{Bases have a unique dual. Redundant frames have infinitely
many.} The freedom is exactly $\ker F$: the more redundant the frame
(the larger $m-n$), the larger the space of valid dual frames. This is
the algebraic shadow of the flexibility we advertised informally all the
way back in Chapter~1, Section~1.1: redundancy buys you freedom of
representation, and Theorem~\ref{thm:dual-family} is the precise
statement of how much freedom, and exactly where it lives.
\end{keyidea}

\begin{examplebox}[title={Example 5.2: An alternate dual of the non-tight frame}]
\label{ex:alternate-dual-h}
Continue Example~\ref{ex:canonical-dual-h}. Since $m=3>n=2$, $\ker F$
is $1$-dimensional. One checks $F(1,1,-1)^T = h_1+h_2-h_3 =
(1,0)^T+(0,1)^T-(1,1)^T=(0,0)^T$, so $(1,1,-1)^T$ spans $\ker F$. A
valid perturbation is $Z = w\,(1,1,-1)$ for any $w\in\R^2$ (a single
column-vector $w$ outer-producted against the row $(1,1,-1)$), since
then each row of $Z$ is a multiple of $(1,1,-1)\in\ker F$, and
$ZF^*=0$ as required. Taking, say, $w=(0.1,-0.2)^T$:
\[
  Z = \begin{pmatrix}0.1\\-0.2\end{pmatrix}\begin{pmatrix}1&1&-1\end{pmatrix}
  = \begin{pmatrix}0.1&0.1&-0.1\\-0.2&-0.2&0.2\end{pmatrix},
\]
and $G = \tilde F + Z$ gives a valid, non-canonical dual frame
$\{g_1,g_2,g_3\}$. This is verified numerically in
Lab Exercise 5.2, including a direct check
that $\sum_i\inner{x}{h_i}g_i=x$ for several test vectors $x$, despite
$\{g_i\}\ne\{\tilde h_i\}$.
\end{examplebox}

\section{The Canonical Dual Is the Minimal-Norm Dual}
\label{sec:minimal-norm}

Among the infinitely many dual frames available whenever $m>n$, why is
the canonical dual the natural default? The next theorem gives a precise
optimality reason.

\begin{theorem}[Minimal-norm characterization of the canonical dual]
\label{thm:canonical-minimal-norm}
Among all dual synthesis matrices $G$ satisfying $GF^*=I_n$, the
canonical dual $\tilde F = S^{-1}F$ is the unique minimizer of the total
squared norm
\[
  \norm{G}_F^2 := \sum_{i=1}^m \norm{g_i}^2 = \tr(GG^*),
\]
(the squared Frobenius norm of $G$). That is, $\tilde F$ uses the
``smallest'' dual vectors, in this total-energy sense, among all valid
choices.
\end{theorem}

\begin{proof}
Write $G = \tilde F + Z$ with $ZF^*=0$, as guaranteed by
Theorem~\ref{thm:dual-family}. We expand the squared Frobenius norm:
\[
  \norm{G}_F^2 = \tr\bigl((\tilde F+Z)(\tilde F+Z)^*\bigr)
  = \tr(\tilde F\tilde F^*) + \tr(\tilde FZ^*) + \tr(Z\tilde F^*) + \tr(ZZ^*).
\]
We claim the two cross terms vanish. Indeed,
\[
  \tilde F Z^* = (S^{-1}F)Z^* = S^{-1}(FZ^*) = S^{-1}(ZF^*)^* = S^{-1}\cdot0^* = 0,
\]
using $ZF^*=0$. Taking adjoints of both sides of $ZF^*=0$ gives
$(ZF^*)^*=FZ^*=0$ as well (using $(AB)^*=B^*A^*$ and $(F^*)^*=F$), which
is exactly the fact used above. Hence $\tr(\tilde FZ^*)=0$, and
similarly $\tr(Z\tilde F^*)=\tr((\tilde FZ^*)^*)=\tr(0^*)=0$. This leaves
\[
  \norm{G}_F^2 = \norm{\tilde F}_F^2 + \norm{Z}_F^2 \;\ge\; \norm{\tilde F}_F^2,
\]
with equality if and only if $\norm{Z}_F^2=0$, i.e.\ $Z=0$, i.e.\
$G=\tilde F$. So $\tilde F$ is the unique minimizer.
\end{proof}

\begin{keyidea}
The identity $\norm{G}_F^2 = \norm{\tilde F}_F^2+\norm{Z}_F^2$ is a
\textbf{Pythagorean decomposition}: the canonical dual's contribution
and the ``extra'' perturbation $Z$ contribute independently
(orthogonally) to the total energy, exactly as $a^2+b^2$ decomposes the
squared length of a right-triangle hypotenuse. This is the same
geometric idea as projecting onto a subspace and decomposing a vector
into a component inside and a component outside: here, the
``subspace'' is the affine family of duals, the canonical dual is the
point of that family closest to the origin, and every other dual is
strictly farther away by exactly $\norm{Z}_F^2$.
\end{keyidea}

\begin{remark*}
This theorem explains, precisely, why the term ``canonical'' is earned
and not arbitrary: among every possible way to reconstruct $x$ from its
frame coefficients $\inner{x}{f_i}$, the canonical dual does so using
the least total ``effort,'' in the sense of minimizing $\sum_i\norm{g_i}^2$.
This will matter in Chapter~11 (frames and quantization/noise), where
smaller dual-vector norms translate directly into better noise
suppression in the reconstruction formula --- a connection we begin to
explore numerically in the next section.
\end{remark*}

\section{Reconstruction Stability: Canonical vs.\ Alternate Duals}
\label{sec:stability}

Chapter~1 motivated frames by the promise of robustness to noise and
erasures. We are now equipped to investigate this quantitatively for the
first time. Suppose the frame coefficients $c_i = \inner{x}{f_i}$ are
corrupted by some noise $\varepsilon_i$, so that the receiver has access
only to $c_i + \varepsilon_i$, and reconstructs via a dual frame
$\{g_i\}$ as $\hat x = \sum_i (c_i+\varepsilon_i) g_i = x +
\sum_i\varepsilon_i g_i$. The reconstruction error is
\[
  \hat x - x = \sum_{i=1}^m \varepsilon_i g_i = G\varepsilon,
\]
where $\varepsilon = (\varepsilon_1,\ldots,\varepsilon_m)$. If the noise
components $\varepsilon_i$ are independent with variance $\sigma^2$
each, the expected squared reconstruction error is
\[
  \mathbb{E}\norm{\hat x-x}^2 = \mathbb{E}\norm{G\varepsilon}^2
  = \sigma^2\,\tr(GG^*) = \sigma^2\norm{G}_F^2.
\]

\begin{keyidea}
\textbf{The canonical dual minimizes expected reconstruction error under
uncorrelated noise.} This follows immediately from
Theorem~\ref{thm:canonical-minimal-norm}: since
$\mathbb{E}\norm{\hat x-x}^2=\sigma^2\norm{G}_F^2$ and $\tilde F$
minimizes $\norm{G}_F^2$ among all valid duals, the canonical dual also
minimizes expected reconstruction error. The abstract optimality result
of Section~\ref{sec:minimal-norm} therefore has an immediate and very
concrete payoff: \emph{when in doubt, use the canonical dual}---it is
not just the most natural choice algebraically, but the most
noise-robust choice as well.
\end{keyidea}

\begin{examplebox}[title={Example 5.3: Numerically comparing canonical vs.\ alternate dual}]
\label{ex:stability-comparison}
Continue Examples~\ref{ex:canonical-dual-h} and~\ref{ex:alternate-dual-h}.
The canonical dual has $\norm{\tilde F}_F^2 = \tr(S^{-1})$. Since
$S^{-1}=\tfrac13\begin{pmatrix}2&-1\\-1&2\end{pmatrix}$, we get
$\tr(S^{-1})=\tfrac13(2+2)=\tfrac43\approx1.333$. For the alternate dual
$G=\tilde F+Z$ from Example~\ref{ex:alternate-dual-h} with
$Z=(0.1,-0.2)^T(1,1,-1)$, each row of $Z$ has \emph{three} entries (one
per frame vector), so direct computation (carried out numerically in
Lab Exercise 5.3) gives
$\norm{Z}_F^2=3(0.1^2+0.2^2)=0.15$, so
$\norm{G}_F^2=\norm{\tilde F}_F^2+\norm{Z}_F^2\approx1.333+0.15=1.483$,
strictly larger than the canonical dual's $1.333$, exactly as
Theorem~\ref{thm:canonical-minimal-norm} guarantees. Under unit-variance
noise, the canonical dual's expected squared error is $\approx1.333$,
while the alternate dual's is $\approx1.483$: roughly $11.25\%$ worse,
purely as a consequence of choosing a non-optimal dual.
\end{examplebox}

\section{NumPy Implementation}
\label{sec:numpy-dual-frames}

\begin{lstlisting}[caption={Canonical and alternate dual frames, and stability comparison}]
import numpy as np
from scipy.linalg import null_space

def frame_operator(F):
    return F @ F.T

def canonical_dual(F):
    """Synthesis matrix of the canonical dual frame: S^{-1} F."""
    S = frame_operator(F)
    return np.linalg.solve(S, F)   # S^{-1} @ F, more stable than explicit inverse

def dual_family_direction(F):
    """
    Orthonormal basis of directions Z = w @ N.T (one column w at a time)
    spanning the freedom in choosing a dual frame, where N spans ker(F).
    Returns N, shape (m, m-n).
    """
    return null_space(F)

def random_alternate_dual(F, scale=1.0, seed=0):
    """Construct a random valid (non-canonical) dual frame's synthesis matrix."""
    n, m = F.shape
    N = dual_family_direction(F)        # shape (m, m-n)
    rng = np.random.default_rng(seed)
    W = rng.standard_normal((n, N.shape[1])) * scale
    Z = W @ N.T                          # shape (n, m), and Z @ F.T == 0
    G_tilde = canonical_dual(F)
    return G_tilde + Z

def expected_reconstruction_error(G, sigma=1.0):
    """E[||x_hat - x||^2] / sigma^2 = ||G||_F^2 = tr(G G^T)."""
    return np.trace(G @ G.T)

# -----------------------------------------------------------------
# Example 5.1 / 5.3: non-tight frame {h1, h2, h3}
# -----------------------------------------------------------------
H = np.array([[1., 0., 1.],
              [0., 1., 1.]])

H_tilde = canonical_dual(H)
print("=== Canonical dual of {h1,h2,h3} ===")
print("H_tilde =\n", H_tilde.round(4))
print("Check H_tilde @ H.T == I:", np.allclose(H_tilde @ H.T, np.eye(2)))
print(f"||H_tilde||_F^2 = {expected_reconstruction_error(H_tilde):.6f}")

G_alt = random_alternate_dual(H, scale=0.2, seed=0)
print("\n=== A random alternate dual ===")
print("G_alt =\n", G_alt.round(4))
print("Check G_alt @ H.T == I:", np.allclose(G_alt @ H.T, np.eye(2)))
print(f"||G_alt||_F^2 = {expected_reconstruction_error(G_alt):.6f}")
print(f"(Canonical is smaller: "
      f"{expected_reconstruction_error(H_tilde) <= expected_reconstruction_error(G_alt)})")

# -----------------------------------------------------------------
# Verify reconstruction works exactly for BOTH duals
# -----------------------------------------------------------------
rng = np.random.default_rng(1)
for trial in range(3):
    x = rng.standard_normal(2)
    c = H.T @ x   # frame coefficients <x, h_i>
    x_hat_canonical = H_tilde @ c
    x_hat_alt = G_alt @ c
    print(f"\nx = {x.round(4)}")
    print(f"  canonical reconstruction error: {np.linalg.norm(x_hat_canonical - x):.2e}")
    print(f"  alternate reconstruction error: {np.linalg.norm(x_hat_alt - x):.2e}")

# -----------------------------------------------------------------
# Monte Carlo: verify canonical dual minimizes noisy reconstruction error
# -----------------------------------------------------------------
def noisy_reconstruction_trial(F, G, x, sigma, rng):
    c = F.T @ x
    noise = rng.standard_normal(c.shape) * sigma
    x_hat = G @ (c + noise)
    return np.linalg.norm(x_hat - x)**2

rng = np.random.default_rng(2)
x_test = np.array([1.0, -0.5])
sigma = 0.3
N_trials = 200_000

errs_canonical = [noisy_reconstruction_trial(H, H_tilde, x_test, sigma, rng)
                   for _ in range(N_trials)]
errs_alt = [noisy_reconstruction_trial(H, G_alt, x_test, sigma, rng)
            for _ in range(N_trials)]

print(f"\n=== Monte Carlo noise comparison ({N_trials} trials, sigma={sigma}) ===")
print(f"Mean squared error, canonical dual: {np.mean(errs_canonical):.6f}  "
      f"(predicted: {sigma**2 * expected_reconstruction_error(H_tilde):.6f})")
print(f"Mean squared error, alternate dual: {np.mean(errs_alt):.6f}  "
      f"(predicted: {sigma**2 * expected_reconstruction_error(G_alt):.6f})")
\end{lstlisting}

\begin{remark*}
We use \texttt{np.linalg.solve(S, F)} rather than explicitly forming
\texttt{np.linalg.inv(S) @ F}. These are mathematically identical, but
\texttt{solve} avoids explicitly constructing the inverse matrix, which
is both faster and numerically more stable, especially as $S$ becomes
ill-conditioned (nearly singular). This is a general good habit: prefer
\texttt{np.linalg.solve(A, b)} over \texttt{np.linalg.inv(A) @ b}
whenever you only need the product, not the inverse matrix itself.
\end{remark*}

\section{Chapter Summary}
\label{sec:summary-ch5}

\begin{itemize}
  \item A \textbf{dual frame} $\{g_i\}$ of $\{f_i\}$ satisfies
        $x=\sum_i\inner{x}{f_i}g_i$ for all $x$, equivalently
        $GF^*=I_n$ for the synthesis matrices $G, F$. The relationship
        is mutual: $\{f_i\}$ is also dual to $\{g_i\}$.
  \item The \textbf{canonical dual} $\tilde f_i=S^{-1}f_i$ is always a
        valid dual, with frame operator $\tilde S=S^{-1}$ and frame
        bounds $\tilde A=1/B$, $\tilde B=1/A$ (reciprocal and reversed).
        For tight frames, $\tilde f_i=\tfrac1Af_i$, a trivial rescaling.
  \item \textbf{Theorem~\ref{thm:dual-family}}: every valid dual has
        synthesis matrix $G=\tilde F+Z$ with $ZF^*=0$, i.e.\ each row
        of $Z$ lies in $\ker F$. Bases ($m=n$) have a unique dual;
        redundant frames ($m>n$) have an $(m-n)$-real-dimensional
        affine family of duals.
  \item \textbf{Theorem~\ref{thm:canonical-minimal-norm}}: the canonical
        dual is the unique \textbf{minimal Frobenius-norm} dual, via the
        Pythagorean decomposition $\norm{G}_F^2=\norm{\tilde F}_F^2+\norm{Z}_F^2$.
  \item This optimality has a direct payoff: since expected squared
        reconstruction error under uncorrelated coefficient noise equals
        $\sigma^2\norm{G}_F^2$, the \textbf{canonical dual minimizes
        expected reconstruction error} among all valid duals---the
        precise, quantitative form of the robustness-to-noise promise
        first made informally in Chapter~1.
\end{itemize}

\section{Lab Exercises}
\label{sec:lab-ch5}

\begin{exercisebox}[title={Lab Exercise 5.1: The Canonical Dual Pipeline}]
\label{lab:dual-symmetry}
Implement \texttt{canonical\_dual(F)} from
Section~\ref{sec:numpy-dual-frames}. Run it on:
\begin{enumerate}[label=(\alph*)]
  \item The triangular frame (tight, Chapter~3 Example~3.2): verify
        $\tilde f_i = \tfrac{1}{1.5}f_i$ for each $i$.
  \item The non-tight frame $\{h_1,h_2,h_3\}$ (Example~5.1): reproduce
        the hand computation exactly.
  \item The square vertices (Chapter~3 Example~3.6): verify the
        canonical dual is again a rescaling of the original (as
        expected, since this frame is tight).
\end{enumerate}
For each case, also verify the mutual-duality claim from the remark
after Definition~5.2.1: check numerically that
\texttt{F @ H\_tilde.T} also equals the identity (not just
\texttt{H\_tilde @ F.T}), confirming $\{f_i\}$ is dual to
$\{\tilde f_i\}$ as well as vice versa.
\end{exercisebox}

\begin{exercisebox}[title={Lab Exercise 5.2: Verifying the Dual Family}]
\label{lab:dual-family-numerics}
Implement \texttt{dual\_family\_direction(F)} and
\texttt{random\_alternate\_dual(F)} from
Section~\ref{sec:numpy-dual-frames}. Using $\{h_1,h_2,h_3\}$:
\begin{enumerate}[label=(\alph*)]
  \item Reproduce Example~5.2's specific alternate dual exactly (using
        $w=(0.1,-0.2)^T$), and verify $GF^*=I_2$ numerically.
  \item Generate $10$ different random alternate duals (varying the
        seed) and verify, for each, that $\sum_i\inner{x}{h_i}g_i=x$
        holds exactly (to machine precision) for five random test
        vectors $x$.
  \item For each of the $10$ alternate duals, compute $\norm{G}_F^2$
        and confirm it always exceeds $\norm{\tilde F}_F^2 \approx
        1.333$ (the canonical value from Example~5.3), consistent with
        Theorem~\ref{thm:canonical-minimal-norm}.
\end{enumerate}
\end{exercisebox}

\begin{exercisebox}[title={Lab Exercise 5.3: Monte Carlo Stability Comparison}]
\label{lab:stability-numerics}
Reproduce the Monte Carlo experiment from
Section~\ref{sec:numpy-dual-frames} (the final code block). Then extend
it:
\begin{enumerate}[label=(\alph*)]
  \item Vary the perturbation scale in
        \texttt{random\_alternate\_dual} over
        $\texttt{scale}=0.05, 0.2, 0.5, 1.0, 2.0$, and for each, compute
        $\norm{G}_F^2$ and the predicted mean squared error
        $\sigma^2\norm{G}_F^2$ (with $\sigma=0.3$ as before). Plot
        predicted MSE versus \texttt{scale}.
  \item Confirm, by running the Monte Carlo simulation (at least
        $50{,}000$ trials) at each scale, that the empirically observed
        mean squared error matches your prediction closely.
  \item In a sentence or two: does this experiment suggest any practical
        guidance for choosing a dual frame in a real signal-processing
        system where the frame is fixed in advance but coefficient noise
        is unavoidable?
\end{enumerate}
\end{exercisebox}

\begin{exercisebox}[title={Lab Exercise 5.4: Dual Frames of a Basis Are Unique}]
Let $\{v_1, v_2\} = \{(2,1)^T, (1,3)^T\}$, a basis (not a redundant
frame) for $\R^2$. Compute the canonical dual $\{\tilde v_1, \tilde
v_2\}$, and verify, using \texttt{scipy.linalg.null\_space}, that
$\ker F$ is trivial (zero-dimensional) for the synthesis matrix of this
basis. Conclude, via Theorem~\ref{thm:dual-family}, that the canonical
dual is the \emph{only} dual frame in this case. (For comparison with
your linear algebra background: explain in a sentence why this matches
the familiar fact that a basis has a unique dual/biorthogonal basis,
characterized by $\inner{v_i}{\tilde v_j} = \delta_{ij}$.)
\end{exercisebox}

\begin{exercisebox}[title={Lab Exercise 5.5 (Optional, Challenge): Sparsity-Promoting Duals}]
Theorem~\ref{thm:canonical-minimal-norm} shows the canonical dual
minimizes the \emph{total Frobenius energy} $\sum_i\norm{g_i}^2$ among
all valid duals. A different, and increasingly important, optimality
criterion in modern signal processing is \emph{sparsity}: finding a dual
frame (or, more precisely, a coefficient sequence) with as many exactly
zero entries as possible. This exercise previews that idea numerically.
\begin{enumerate}[label=(\alph*)]
  \item Using $\{h_1,h_2,h_3\}$ and a fixed $x=(1,-0.5)^T$, the
        canonical dual gives one specific coefficient-free
        reconstruction. But recall from
        Example~3.7 that the coefficients $c_i=\inner{x}{f_i}$
        \emph{themselves} are not unique either, once you allow
        $c'=c+z$ for $z\in\ker F$ (a fact distinct from, but related
        to, the dual-frame freedom of this chapter). Parametrize
        $c'(t) = c + t\cdot n$ where $n$ spans $\ker H$ (one-dimensional,
        as in Example~5.2), and plot $\norm{c'(t)}_1$ (the $\ell^1$ norm,
        \texttt{np.sum(np.abs(...))}) as a function of $t\in[-2,2]$.
  \item Find the value of $t$ that minimizes $\norm{c'(t)}_1$
        numerically (e.g.\ via \texttt{scipy.optimize.minimize\_scalar}).
        Is the minimizing coefficient vector sparse (does it have an
        exact zero entry)? This phenomenon---$\ell^1$ minimization
        promoting exact sparsity---is the starting point of compressed
        sensing, which we will touch on again in Chapter~12.
\end{enumerate}
\end{exercisebox}
\newpage
\section{Supplementary Exercises}
\label{sec:extra-ch5}

\subsection*{Theoretical Exercises}
\begin{theoexercises}
	\item Prove that the set of all dual frames of a fixed frame
	$\{f_i\}_{i=1}^m$ (the affine family of Theorem~5.4.1) is
	\emph{convex}: if $G_1$ and $G_2$ are both valid dual synthesis
	matrices, then so is $tG_1+(1-t)G_2$ for every $t\in[0,1]$.
	
	\item Prove the full biconditional: a frame $\{f_i\}_{i=1}^m$ equals its
	own canonical dual ($\tilde f_i=f_i$ for all $i$) if and only if
	$\{f_i\}$ is a Parseval frame. (One direction is immediate;
	for the converse, show $F=\tilde F$ forces $S=I$.)
	
	\item Prove that $\norm{Z}_F^2$ is a strictly convex function of $Z$
	(over the subspace of matrices with $ZF^*=0$), and use this
	convexity --- rather than the Pythagorean/orthogonality argument
	used in the text --- to give an alternative proof that the
	canonical dual is the \emph{unique} minimizer of total energy in
	Theorem~5.5.1.
	
	\item Let $\{f_i\}_{i=1}^m$ be a tight frame for $\R^n$ with bound $A$
	(not necessarily unit-norm). Prove that the canonical dual has
	total energy $\sum_i\norm{\tilde f_i}^2=n/A$ exactly. Specialize
	to the unit-norm case ($A=m/n$) to recover the formula $n^2/m$
	used in Chapter~11.
	
	\item Prove directly from the definition of the analysis operator
	$F^*$ that $\ker F^*=(\spn\{f_i\})^\perp$.
	
	\item Using strict convexity (T3), prove or disprove: for any two
	\emph{distinct} valid duals $G_1\ne G_2$ of a fixed frame, the
	averaged dual $\tfrac12(G_1+G_2)$ has strictly smaller total
	energy than the larger of $\norm{G_1}_F^2,\norm{G_2}_F^2$.
\end{theoexercises}

\subsection*{Computational Exercises}
\begin{compexercises}
	\item For $\{h_1,h_2,h_3\}$, construct two distinct alternate duals
	$G_1,G_2$, form their average, and confirm numerically that the
	average is itself a valid dual (T1).
	
	\item Confirm numerically that the triangular frame (tight, $A=1.5\ne1$)
	does \emph{not} equal its own canonical dual, while the standard
	ONB (Parseval) does, illustrating both directions of T2.
	
	\item Scale the triangular frame by a factor $c=2,3,4$ (multiply every
	vector by $c$), and for each, compute the canonical dual's total
	energy and compare to the predicted $n/A$ from T4.
	
	\item For $\{h_1,h_2,h_3\}$, construct $10$ random alternate duals with
	varying perturbation scale, compute each one's total energy, and
	verify numerically that the canonical dual ($w=0$) has strictly
	smaller energy than every one of them, consistent with T3's
	convexity argument.
	
	\item Construct three distinct alternate duals of $\{h_1,h_2,h_3\}$,
	compute all pairwise averages, and verify each pairwise average
	has energy strictly less than the larger of its two parents,
	illustrating T6 across multiple pairs.
\end{compexercises}